\documentclass{article}
\input{../../../lib.sty}

\title{\texttt{fn solve\_nb\_x\_from\_mean}}
\author{Michael Shoemate}

\begin{document}
\maketitle

Proves soundness of \texttt{fn solve\_nb\_x\_from\_mean} in \texttt{mod.rs}.

\section{Hoare Triple}
\subsection*{Precondition}
\texttt{mean} $> 1$ and \texttt{eta} $\ge 0$.

\subsection*{Pseudocode}
\lstinputlisting[language=Python,firstline=2,escapechar=|]{./pseudocode/solve_nb_x_from_mean.py}

\subsection*{Postcondition}
\begin{theorem}
\label{postcondition}
For any inputs satisfying the precondition, \texttt{solve\_nb\_x\_from\_mean} either returns
\texttt{Err(e)}, or returns \texttt{Ok(out)} where
\[
\texttt{expected\_nb\_mean\_from\_x\_hi(eta, out)} \le \texttt{mean}.
\]
Moreover, \texttt{out} is the final value of the upper bisection endpoint \texttt{hi}.
\end{theorem}

\section{Proof}
Assume the precondition.

\begin{lemma}
\label{err-e}
\texttt{solve\_nb\_x\_from\_mean} only returns \texttt{Err(e)} if some call to
\texttt{expected\_nb\_mean\_from\_x\_hi} propagates \texttt{Err(e)} or if the explicit finite
guard on line \ref{line:finite_guard} fails.
\end{lemma}

\begin{proof}
These are the only fallible operations in the pseudocode.
\end{proof}

\begin{lemma}
\label{hi-invariant}
Whenever control reaches the top of the bisection loop on line \ref{line:bisection},
\[
\texttt{expected\_nb\_mean\_from\_x\_hi(eta, hi)} \le \texttt{mean}.
\]
\end{lemma}

\begin{proof}
When the expansion loop terminates normally, the guard on line \ref{line:expand} is false, so the
current \texttt{hi} satisfies the inequality.
Inside the bisection loop, \texttt{hi} is updated only in the \texttt{else} branch, and there it is
set to \texttt{mid} only after checking that
\texttt{expected\_nb\_mean\_from\_x\_hi(eta, mid)} $\le$ \texttt{mean}.
Thus the property is preserved from one iteration to the next.
\end{proof}

\begin{lemma}
\label{termination}
If the function does not return \texttt{Err(e)}, then the bisection loop eventually reaches line
\ref{line:converged} and returns.
\end{lemma}

\begin{proof}
At the start of the bisection loop, \texttt{lo} and \texttt{hi} are finite and satisfy
\texttt{lo < hi}. Each non-terminating iteration sets exactly one endpoint to the midpoint on line
\ref{line:mid}, so the new interval is a strict subinterval of the previous one and still has
finite endpoints.

Because \texttt{f64} has only finitely many representable values in any bounded interval, this
strict shrinking cannot continue indefinitely. Therefore after finitely many iterations there is no
representable \texttt{f64} strictly between \texttt{lo} and \texttt{hi}. At that point the
midpoint rounds to one of the endpoints, so the test on line \ref{line:converged} succeeds and the
function returns.
\end{proof}

\begin{lemma}
\label{ok-out}
If \texttt{solve\_nb\_x\_from\_mean} returns \texttt{Ok(out)}, then
\[
\texttt{expected\_nb\_mean\_from\_x\_hi(eta, out)} \le \texttt{mean},
\]
and \texttt{out} is the final value of \texttt{hi}.
\end{lemma}

\begin{proof}
The only successful return is on line \ref{line:converged}, which returns \texttt{hi}; therefore
\texttt{out} is exactly the final value of the upper endpoint.

By Lemma~\ref{hi-invariant}, the displayed inequality holds at the start of every bisection
iteration. The convergence test does not modify \texttt{hi}, so it still holds at return.
\end{proof}

\begin{proof}[Proof of Theorem~\ref{postcondition}]
Holds by Lemma~\ref{err-e}, Lemma~\ref{termination}, and Lemma~\ref{ok-out}.
\end{proof}

\end{document}
